\section{Scientific and Practical Significance} \label{section:scientific_and_practical_significance}

Material interfaces are central to the behaviour of functional solids, especially at the nanoscale where surface and interfacial effects dominate bulk properties. In semiconductors, they define charge transport pathways by setting potential barriers and band offsets. In batteries, they regulate ionic conductivity, chemical stability, and degradation kinetics. In photovoltaics and quantum devices, interfacial dipoles, tunnelling channels, and confinement effects underpin key device operations.

Subtle differences in stacking, coherence, or atomic reconstruction can shift critical electronic properties by several electronvolts; altering conductivity, carrier lifetimes, or activation barriers. These shifts emerge from a complex interplay of interfacial strain, chemical bonding, and electronic redistribution, and are sensitive to both global structure and local environment. As such, predicting interface stability and functionality from atomic configuration alone remains a demanding task.

While Density Functional Theory (DFT) offers a robust ab initio framework for computing interfacial energies and electronic properties, its computational cost scales poorly with system size. Realistic supercells, often comprising hundreds of atoms, are therefore difficult to study in large numbers. DFT is well-suited to final validation, but ill-suited to exhaustive exploration across configurational space.

Machine learning (ML) models, particularly machine-learned interatomic potentials (MLPs), provide a way forward. When trained on DFT data, MLPs can reproduce formation energies and atomic forces at orders-of-magnitude lower cost, making them suitable for broad screening applications.  Their integration into automated workflows, using tools such as ARTEMIS and RAFFLE, enables systematic generation and evaluation of interface candidates across crystallographic orientations, stacking shifts, and chemical combinations.

Importantly, the utility of ML does not stop at static energetics. Many applications require insight into how interfaces behave under external stimuli; electric fields, temperature gradients, or chemical exposure. More adaptive and scalable modelling approaches are required for emerging interfacial phases, such as field-tunable dielectric layers or metastable heterostructures.

This project contributes to that broader effort. It assesses the ability of MLPs, specifically MACE, to recover DFT-derived favourability rankings, identifies where discrepancies arise, and proposes delta-learning corrections to address them. It also lays groundwork for predictive models that estimate interface properties from top and bottom slabs, enabling pre-screening before full interface construction. By focusing on group-IV semiconductor systems and both homostructure and heterostructure cases, this work targets a class of technologically relevant interfaces where modelling uncertainty remains high and conventional rules often fail.

Ultimately, this research supports the advancement of data-driven, scalable methods for interface discovery; bridging the gap between computational feasibility and the physical fidelity required for next-generation materials design.